
\documentclass[12pt,a4paper,openany,oneside]{book}

% --- Paquetes ---
\usepackage[utf8]{inputenc}
\usepackage[spanish, es-noshorthands]{babel}
\usepackage{amsmath, amsfonts, amssymb}
\usepackage{graphicx}
\usepackage{geometry}
\usepackage{hyperref}
\usepackage{booktabs}
\usepackage{fancyhdr}
\usepackage{listings}
\usepackage{xcolor}
\usepackage{tikz}
\usepackage{pgfplots}
\usetikzlibrary{arrows.meta, positioning, shapes.geometric}
\pgfplotsset{compat=1.18}

\geometry{margin=2.5cm}

% --- Colores y estilos para código Python ---
\definecolor{codegreen}{rgb}{0,0.6,0}
\definecolor{codegray}{rgb}{0.5,0.5,0.5}
\definecolor{codepurple}{rgb}{0.58,0,0.82}
\definecolor{backcolour}{rgb}{0.95,0.95,0.92}

\lstset{
    language=Python,
    backgroundcolor=\color{backcolour},   
    commentstyle=\color{codegreen},
    keywordstyle=\color{blue},
    numberstyle=\tiny\color{codegray},
    stringstyle=\color{codepurple},
    basicstyle=\ttfamily\small,
    breaklines=true,
    frame=single,
    showstringspaces=false
}

% --- Estilo de cabecera y pie ---
\pagestyle{fancy}
\fancyhf{}
\renewcommand{\headrulewidth}{0.4pt}
\renewcommand{\footrulewidth}{0.4pt}
\fancyhead[L]{\small Carlos César Sánchez Coronel}
\fancyhead[R]{\small Sesión 13: Descomposición en Valores Singulares (SVD)}
\fancyfoot[L]{\small Carlos César Sánchez Coronel}
\fancyfoot[C]{\thepage}

% --- Portada ---
\title{
    \textbf{Sesión 13} \\ 
    \Huge \textbf{DESCOMPOSICIÓN EN VALORES SINGULARES (SVD)} \\
    \large La navaja suiza del álgebra lineal \\
    \normalsize Aplicaciones en compresión de imágenes, sistemas de recomendación, LSA y más
}
\author{
    \small Por \\ \vspace{0.2cm}
    \Large \textsc{Carlos César Sánchez Coronel}
}
\date{2026}

\begin{document}

\maketitle
\tableofcontents
\addcontentsline{toc}{chapter}{Índice general}

% ------------------------------------------------------------------
% CAPÍTULO 13: SESIÓN 13 - SVD
% ------------------------------------------------------------------
\chapter{Descomposición en Valores Singulares (SVD)}

La descomposición en valores singulares (SVD, por sus siglas en inglés) es una de las herramientas más poderosas y versátiles del álgebra lineal. Nos permite factorizar cualquier matriz (no necesariamente cuadrada) en el producto de tres matrices con una interpretación geométrica clara: rotaciones y escalados. Sus aplicaciones en inteligencia artificial son innumerables: compresión de imágenes, reducción de ruido, sistemas de recomendación, análisis semántico latente (LSA), cálculo de pseudoinversa, inicialización de redes neuronales, y mucho más. En esta sesión exploraremos en profundidad el teorema SVD, su interpretación, y cómo implementarlo en Python para resolver problemas reales.

\section{Teorema de la descomposición en valores singulares}

\subsection{Enunciado del teorema}
Sea $\mathbf{A} \in \mathbb{R}^{m \times n}$ una matriz real de rango $r$. Entonces existen matrices ortogonales $\mathbf{U} \in \mathbb{R}^{m \times m}$, $\mathbf{V} \in \mathbb{R}^{n \times n}$ y una matriz diagonal $\boldsymbol{\Sigma} \in \mathbb{R}^{m \times n}$ tales que:
\[
\mathbf{A} = \mathbf{U} \boldsymbol{\Sigma} \mathbf{V}^\top
\]
donde:
\begin{itemize}
    \item $\mathbf{U} = [\mathbf{u}_1 \dots \mathbf{u}_m]$ contiene los \textbf{vectores singulares izquierdos} (ortonormales).
    \item $\mathbf{V} = [\mathbf{v}_1 \dots \mathbf{v}_n]$ contiene los \textbf{vectores singulares derechos} (ortonormales).
    \item $\boldsymbol{\Sigma}$ tiene la forma:
    \[
    \boldsymbol{\Sigma} = \begin{pmatrix}
    \sigma_1 & 0 & \cdots & 0 & \cdots & 0 \\
    0 & \sigma_2 & \cdots & 0 & \cdots & 0 \\
    \vdots & \vdots & \ddots & \vdots & & \vdots \\
    0 & 0 & \cdots & \sigma_r & \cdots & 0 \\
    \vdots & \vdots & & \vdots & \ddots & \vdots \\
    0 & 0 & \cdots & 0 & \cdots & 0
    \end{pmatrix}
    \]
    con $\sigma_1 \ge \sigma_2 \ge \dots \ge \sigma_r > 0$ llamados \textbf{valores singulares}.
\end{itemize}
Las columnas de $\mathbf{U}$ son autovectores de $\mathbf{A}\mathbf{A}^\top$, las columnas de $\mathbf{V}$ son autovectores de $\mathbf{A}^\top\mathbf{A}$, y los valores singulares son las raíces cuadradas de los autovalores comunes (no nulos) de esas matrices.

\subsection{Interpretación geométrica}
La SVD descompone la transformación lineal representada por $\mathbf{A}$ en tres pasos:
\begin{enumerate}
    \item Una rotación (o reflexión) en el espacio de entrada dada por $\mathbf{V}^\top$.
    \item Un escalado a lo largo de los ejes coordenados por los valores singulares (y posiblemente un cambio de dimensión) mediante $\boldsymbol{\Sigma}$.
    \item Una rotación (o reflexión) en el espacio de salida dada por $\mathbf{U}$.
\end{enumerate}

\begin{center}
\begin{tikzpicture}
\draw[->] (-1,0) -- (3,0) node[right] {$\mathbb{R}^n$};
\draw[->] (0,-1) -- (0,3) node[above] {$\mathbb{R}^m$};
\draw[blue, thick, ->] (0,0) -- (1,0.5) node[midway, above] {$\mathbf{v}_1$};
\draw[red, thick, ->] (0,0) -- (0.2,1) node[midway, left] {$\mathbf{v}_2$};
\node at (0.5,1.8) {$\mathbf{A}$};
\draw[->, thick] (2,1) -- (4,2) node[midway, above] {$\mathbf{U}$};
\end{tikzpicture}
\captionof{figure}{Interpretación geométrica: $\mathbf{A} = \mathbf{U}\boldsymbol{\Sigma}\mathbf{V}^\top$.}
\end{center}

\subsection{Forma reducida (economy SVD)}
Cuando $m > n$ (más filas que columnas), podemos considerar la versión reducida:
\[
\mathbf{A} = \mathbf{U}_r \boldsymbol{\Sigma}_r \mathbf{V}^\top
\]
donde $\mathbf{U}_r \in \mathbb{R}^{m \times r}$ contiene solo los primeros $r$ vectores singulares izquierdos, $\boldsymbol{\Sigma}_r \in \mathbb{R}^{r \times r}$ es la diagonal con los $r$ valores singulares, y $\mathbf{V} \in \mathbb{R}^{n \times r}$ (o $\mathbf{V}_r$). Esta versión ahorra memoria y cómputo.

\section{Propiedades importantes}

\begin{itemize}
    \item Los valores singulares son únicos (salvo el orden).
    \item $\|\mathbf{A}\|_2 = \sigma_1$ (norma espectral).
    \item La norma de Frobenius: $\|\mathbf{A}\|_F = \sqrt{\sum_{i=1}^r \sigma_i^2}$.
    \item El rango de $\mathbf{A}$ es el número de valores singulares no nulos.
    \item Las columnas de $\mathbf{U}$ correspondientes a valores singulares no nulos forman una base ortonormal del espacio columna de $\mathbf{A}$; las de $\mathbf{V}$ forman una base ortonormal del espacio fila.
    \item Las columnas de $\mathbf{U}$ correspondientes a valores singulares nulos forman una base del complemento ortogonal (espacio nulo izquierdo), y las de $\mathbf{V}$ correspondientes a valores singulares nulos forman una base del espacio nulo.
\end{itemize}

\section{Aproximación de bajo rango y compresión}

Una de las aplicaciones más importantes de la SVD es la aproximación de una matriz por otra de rango inferior. El teorema de Eckart-Young-Mirsky establece que la mejor aproximación de rango $k$ a $\mathbf{A}$ (en norma de Frobenius o espectral) se obtiene truncando la SVD a los primeros $k$ valores singulares:
\[
\mathbf{A}_k = \mathbf{U}_k \boldsymbol{\Sigma}_k \mathbf{V}_k^\top = \sum_{i=1}^k \sigma_i \mathbf{u}_i \mathbf{v}_i^\top
\]
Esto se utiliza en compresión de imágenes: una imagen es una matriz de píxeles; podemos guardar solo los primeros $k$ valores singulares y sus correspondientes vectores, reduciendo drásticamente el espacio de almacenamiento con una pérdida controlada de calidad.

\section{Pseudoinversa de Moore-Penrose}

Para una matriz $\mathbf{A}$ (no necesariamente cuadrada ni de rango completo), la pseudoinversa $\mathbf{A}^+$ generaliza la noción de inversa. Una definición usando SVD es:
\[
\mathbf{A}^+ = \mathbf{V} \boldsymbol{\Sigma}^+ \mathbf{U}^\top
\]
donde $\boldsymbol{\Sigma}^+$ se obtiene tomando el recíproco de cada valor singular no nulo y dejando ceros en los demás:
\[
\boldsymbol{\Sigma}^+ = \operatorname{diag}\left(\frac{1}{\sigma_1}, \dots, \frac{1}{\sigma_r}, 0, \dots, 0\right)
\]
La pseudoinversa tiene la propiedad de que $\mathbf{x} = \mathbf{A}^+ \mathbf{b}$ es la solución de mínimos cuadrados de norma mínima del sistema $\mathbf{A}\mathbf{x} \approx \mathbf{b}$. Es decir, minimiza $\|\mathbf{A}\mathbf{x} - \mathbf{b}\|_2$ y, entre todas las soluciones que lo hacen, elige la de menor norma.

\section{Aplicaciones en Inteligencia Artificial}

\subsection{Compresión de imágenes (Visión por Computador)}
Cada imagen en escala de grises puede verse como una matriz $\mathbf{A} \in \mathbb{R}^{h \times w}$. Aplicamos SVD y truncamos a $k$ valores singulares. La imagen comprimida se reconstruye como $\mathbf{A}_k$. Cuanto menor $k$, mayor compresión pero más pérdida de calidad.

\begin{center}
\begin{tikzpicture}
\draw[fill=gray!30] (0,0) rectangle (4,4);
\node at (2,2) {Imagen original};
\draw[->, thick] (4.5,2) -- (5.5,2);
\draw[fill=gray!30] (6,0) rectangle (8,4);
\node at (7,2) {$\mathbf{U}_k\boldsymbol{\Sigma}_k\mathbf{V}_k^\top$};
\end{tikzpicture}
\captionof{figure}{Compresión: se almacenan solo $k(h+w+1)$ números en lugar de $h\times w$.}
\end{center}

\subsection{Sistemas de recomendación (filtrado colaborativo)}
Supongamos una matriz $\mathbf{R}$ de $m$ usuarios y $n$ productos, donde $R_{ij}$ es la calificación que el usuario $i$ dio al producto $j$ (puede tener valores faltantes). Podemos aproximar $\mathbf{R}$ mediante una factorización de bajo rango $\mathbf{R} \approx \mathbf{U} \mathbf{V}^\top$, con $\mathbf{U} \in \mathbb{R}^{m \times k}$, $\mathbf{V} \in \mathbb{R}^{n \times k}$. Esto equivale a encontrar la mejor aproximación de rango $k$ en el sentido de mínimos cuadrados sobre los valores observados. La SVD de la matriz completa (rellenando con ceros o mediante técnicas iterativas) proporciona esta factorización. Fue la base del algoritmo ganador del Netflix Prize.

\subsection{Análisis Semántico Latente (LSA) en NLP}
Dada una matriz de términos-documentos $\mathbf{A} \in \mathbb{R}^{t \times d}$ (cada columna es un documento, cada fila una palabra, con pesos tf-idf), la SVD nos permite encontrar ``conceptos'' latentes. La aproximación de bajo rango $\mathbf{A}_k$ captura las relaciones semánticas principales, agrupando palabras y documentos que comparten temas. Luego se puede usar para recuperación de información, clustering de documentos, etc.

\subsection{Deep Learning: inicialización y pseudoinversa}
\begin{itemize}
    \item Inicialización de pesos: algunas técnicas utilizan SVD para inicializar capas de manera que se preserve la norma de los gradientes (por ejemplo, inicialización ortogonal).
    \item Pseudoinversa en redes: en algunas arquitecturas (como extreme learning machines) se usa la pseudoinversa para calcular los pesos de salida de forma analítica.
    \item Compresión de modelos: se puede aproximar una capa fully connected de gran tamaño mediante su descomposición SVD de bajo rango, reduciendo el número de parámetros y acelerando la inferencia.
\end{itemize}

\section{Python: SVD con NumPy y aplicaciones}

\subsection{Cálculo de SVD}
\begin{lstlisting}[language=Python]
import numpy as np

# Matriz de ejemplo
A = np.array([[1, 2, 3],
              [4, 5, 6],
              [7, 8, 9]], dtype=float)

U, s, Vt = np.linalg.svd(A, full_matrices=False)
print("U:\n", U)
print("Valores singulares:", s)
print("V^T:\n", Vt)

# Reconstrucción
Sigma = np.diag(s)
A_rec = U @ Sigma @ Vt
print("Reconstrucción (error):", np.linalg.norm(A - A_rec))
\end{lstlisting}

\subsection{Compresión de imágenes}
\begin{lstlisting}[language=Python]
import matplotlib.pyplot as plt
from skimage import data

# Cargar una imagen en escala de grises
img = data.camera().astype(float) / 255.0
print("Forma de la imagen:", img.shape)

# SVD
U, s, Vt = np.linalg.svd(img, full_matrices=False)

# Reconstruir con diferentes rangos
k_values = [5, 20, 50, 100]
plt.figure(figsize=(12,6))
for i, k in enumerate(k_values):
    # Truncar
    Uk = U[:, :k]
    sk = s[:k]
    Vtk = Vt[:k, :]
    img_k = Uk @ np.diag(sk) @ Vtk
    plt.subplot(2, len(k_values), i+1)
    plt.imshow(img_k, cmap='gray')
    plt.title(f'k = {k}')
    plt.axis('off')
    plt.subplot(2, len(k_values), i+1+len(k_values))
    # Error relativo
    error = np.linalg.norm(img - img_k) / np.linalg.norm(img)
    plt.bar(i, error)
    plt.xticks([])
    plt.ylabel('Error relativo')

plt.tight_layout()
plt.show()
\end{lstlisting}

\subsection{Pseudoinversa y regresión lineal}
Resolvemos un sistema sobredeterminado $\mathbf{A}\mathbf{x} \approx \mathbf{b}$ usando la pseudoinversa y comparamos con \texttt{np.linalg.lstsq}.

\begin{lstlisting}[language=Python]
# Datos sintéticos
np.random.seed(42)
m, n = 20, 5
A = np.random.randn(m, n)
x_true = np.random.randn(n)
b = A @ x_true + 0.1 * np.random.randn(m)  # ruido

# Pseudoinversa usando SVD
U, s, Vt = np.linalg.svd(A, full_matrices=False)
# Inversa de los valores singulares
s_inv = np.diag(1.0 / s)
x_pinv = Vt.T @ s_inv @ U.T @ b

# Usando lstsq
x_lstsq, residuos, rango, s_lstsq = np.linalg.lstsq(A, b, rcond=None)

print("Error pseudoinversa vs verdadero:", np.linalg.norm(x_pinv - x_true))
print("Error lstsq vs verdadero:", np.linalg.norm(x_lstsq - x_true))
\end{lstlisting}

\subsection{Simulación de sistema de recomendación}
Generamos una matriz de calificaciones con estructura de bajo rango más ruido y la aproximamos con SVD.

\begin{lstlisting}[language=Python]
# Parámetros
n_users, n_items, k_true = 100, 80, 5
np.random.seed(42)

# Factores latentes verdaderos
U_true = np.random.randn(n_users, k_true)
V_true = np.random.randn(n_items, k_true)
R = U_true @ V_true.T + 0.5 * np.random.randn(n_users, n_items)

# Ocultar algunos valores (simulando datos faltantes)
mask = np.random.rand(n_users, n_items) > 0.3
R_obs = R * mask

# Para completar, podríamos usar SVD sobre la matriz observada rellenando con cero, pero en la práctica se usan métodos iterativos.
# Aquí simplemente aplicamos SVD a R_obs (los ceros afectan)
U, s, Vt = np.linalg.svd(R_obs, full_matrices=False)
k = 5
R_pred = U[:, :k] @ np.diag(s[:k]) @ Vt[:k, :]

# Error en los no observados (test)
test_mask = ~mask
error_test = np.sqrt(np.mean((R[test_mask] - R_pred[test_mask])**2))
print("RMSE en test:", error_test)
\end{lstlisting}

\section{Notación en papers científicos}

En la literatura, la SVD aparece con frecuencia:

\begin{itemize}
    \item En artículos de sistemas de recomendación: $\mathbf{R} \approx \mathbf{U}\mathbf{V}^\top$, donde $\mathbf{U}$ y $\mathbf{V}$ son matrices de factores latentes.
    \item En procesamiento de señales e imágenes: la descomposición en valores singulares de una matriz de datos.
    \item En aprendizaje automático teórico: uso de la pseudoinversa para derivar soluciones de mínimos cuadrados.
    \item En reducción de dimensionalidad: truncamiento de la SVD para obtener una representación de bajo rango.
\end{itemize}

Ejemplo típico:
\[
\mathbf{X} = \mathbf{U} \boldsymbol{\Sigma} \mathbf{V}^\top, \quad \mathbf{X}_k = \sum_{i=1}^k \sigma_i \mathbf{u}_i \mathbf{v}_i^\top
\]

\section{Ejercicios propuestos}

\begin{enumerate}
    \item Calcula la SVD de la matriz $\begin{pmatrix} 1 & 0 \\ 0 & 2 \\ 0 & 0 \end{pmatrix}$ a mano (puedes usar papel y lápiz para encontrar $U$, $\Sigma$, $V$).
    \item Demuestra que la pseudoinversa de $\mathbf{A}$ satisface $\mathbf{A}^+ \mathbf{A} = \mathbf{V}_r \mathbf{V}_r^\top$ (proyección sobre el espacio fila).
    \item Usa SVD para comprimir una imagen a color (tres canales). Para cada canal aplica SVD y visualiza el resultado con distintos $k$.
    \item Implementa un sistema de recomendación simple usando SVD sobre la matriz MovieLens (puedes descargar un subset). Divide en entrenamiento y prueba (ocultando algunos valores) y evalúa el error RMSE.
    \item Investiga el algoritmo de ``randomized SVD'' y cuándo es preferible usarlo en lugar de la SVD clásica.
\end{enumerate}

\section{Resumen}

En esta sesión hemos cubierto:
\begin{itemize}
    \item El teorema de la descomposición en valores singulares y su interpretación geométrica.
    \item La aproximación de bajo rango y su uso en compresión.
    \item La pseudoinversa de Moore-Penrose y su relación con mínimos cuadrados.
    \item Aplicaciones clave en visión (compresión de imágenes), sistemas de recomendación (filtrado colaborativo), NLP (LSA) y deep learning.
    \item Implementaciones prácticas en Python con ejemplos de código.
    \item La notación típica en artículos científicos.
\end{itemize}
La SVD es verdaderamente una navaja suiza: con ella podemos entender la estructura de los datos, reducir dimensionalidad, eliminar ruido y resolver sistemas lineales de manera óptima.

\end{document}
```